\section{Introduction}
\label{sec:introduction}

Digital twins mirror the state of physical assets through continuous data exchange. Cognitive
digital twins go further, adding capabilities such as perception, memory, reasoning and
autonomous decision support, typically built on ontologies and knowledge graphs~\citep{zheng2022cdt}.

\todo{Motivation: predictive maintenance, cost of unplanned downtime.}

\todo{Gap: conceptual / proprietary CDTs; LLM twins without evaluation; no open benchmark.}

\paragraph{Research questions.}
\begin{itemize}
  \item \textbf{RQ1:} Can a modular CDT be built end-to-end from open components and evaluated on a public benchmark?
  \item \textbf{RQ2:} Does grounding an LLM agent in the twin improve answer accuracy and faithfulness?
  \item \textbf{RQ3:} Does knowledge-graph context improve fault localisation and explanation?
  \item \textbf{RQ4:} Does uncertainty-aware prognosis reduce maintenance cost?
\end{itemize}

\paragraph{Contributions.}
\begin{enumerate}
  \item A CDT reference architecture mapped to concrete, swappable components.
  \item \emph{mindtwin}, an open-source, reproducible implementation.
  \item \emph{TwinQA}, a benchmark for evaluating twin-grounded agents.
  \item An empirical study answering RQ1--RQ4.
\end{enumerate}
